\section{Discussion}
\label{sec:discussion}

\subsection{Mean Wiggle Rate is Itself Jagged Across Pressure Levels}
\label{sec:discussion-jagged}

If we define a judge's \textbf{jaggedness} as the standard deviation of its mean wiggle rates across all datasets, 
we can plot each judge's mean wiggle against its jaggedness, separately at each pressure level 
(Figure~\ref{fig:wiggliness-jaggedness}).

\begin{figure}[t]
  \centering
  \includegraphics[width=\linewidth]{per_level_scatter_combined.pdf}
  \caption{Each panel plots all 9 judges at one pressure level: x-axis is the judge's mean wiggle rate at that 
  level (averaged across the 14 judging tasks); y-axis is its cross-dataset jaggedness (std of wiggle rates). 
  Dashed line is the OLS fit, $R^2$ and Pearson $r$ annotated.}
  \label{fig:wiggliness-jaggedness}
\end{figure}

At L1--L5, mean wiggle and cross-dataset jaggedness have a positive linear relationship with fairly strong linear regression fit
($R^2 = 0.68$, $0.94$, $0.87$, $0.58$, $0.75$).
For L1-L5, judges that wiggle more on average also have a wider spread of wiggle rates across datasets. 
At L6, however, the relationship \emph{inverts} ($R^2 = 0.64$, $r = -0.80$): 
judges with the lowest mean wiggle rates (e.g. Gemini 3 Flash, Grok-4.1 R, Gemini 3.1 Pro) have the biggest 
cross-dataset spread. Jaggedness itself takes different shapes across different kinds of pressure.

\subsection{Baseline Jury Majority Strength is a Simple Reliability Screen}
\label{sec:discussion-jury}

Can a cheap test predict which items are likely to be epistemically unstable? We compare three candidate predictors 
of per-item wiggle: \emph{jury majority strength} (size of the L0 
majority across the 9 judges, no pressure applied~\citep{zhao2024council}),
\emph{repeat consistency} (temperature-zero per-item agreement), 
and \emph{position invariance} (verdict survival under argument reordering). 
The strongest predictive signal at every level is \emph{jury majority strength} (Table~\ref{tab:predictor-comparison}, 
mean $|\rho| = 0.59$ vs.\ 0.42 for repeat and 0.37 for invariance). All 84 (dataset, rubric, scale, level)
correlations are negative, with median $|\rho| = 0.58$ (Table~\ref{tab:jury-rho}).

\begin{table}[t]
\centering
\small
\caption{Mean $|\rho|$ between each predictor and per-item wiggle rate, by level. Jury averaged over 84 
(dataset, rubric, scale, level) cells; Repeat and Invariance over 72 (not measured on WildGuard). 
Per-cell jury correlations are reported in Table~\ref{tab:jury-rho}.}
\label{tab:predictor-comparison}
\begin{tabular}{@{}lrrrrrrr@{}}
\toprule
\textbf{Predictor} & \textbf{L1} & \textbf{L2} & \textbf{L3} & \textbf{L4} & \textbf{L5} & \textbf{L6} & \textbf{Overall} \\
\midrule
Jury Majority Strength & \textbf{0.65} & \textbf{0.57} & \textbf{0.57} & \textbf{0.60} & \textbf{0.59} & \textbf{0.57} & \textbf{0.59} \\
Repeat (temp=0)        & 0.44 & 0.38 & 0.39 & 0.42 & 0.42 & 0.44 & 0.42 \\
Invariance (position)  & 0.35 & 0.36 & 0.38 & 0.35 & 0.37 & 0.38 & 0.37 \\
\bottomrule
\end{tabular}
\end{table}

\begin{table}[t]
\centering
\small
\caption{Spearman $\rho$ between baseline jury majority strength and per-item wiggle rate, for each (dataset, rubric, scale, level) cell. All 84 cells are negative.}
\label{tab:jury-rho}
\begin{tabular}{@{}llrrrrrr@{}}
\toprule
\textbf{Dataset} & \textbf{Scale} & \textbf{L1} & \textbf{L2} & \textbf{L3} & \textbf{L4} & \textbf{L5} & \textbf{L6} \\
\midrule
WildGuard      & binary & $-0.684$ & $-0.683$ & $-0.678$ & $-0.623$ & $-0.675$ & $-0.653$ \\
WildGuard      & Likert & $-0.640$ & $-0.379$ & $-0.437$ & $-0.708$ & $-0.766$ & $-0.683$ \\
AEGIS          & binary & $-0.661$ & $-0.628$ & $-0.600$ & $-0.706$ & $-0.655$ & $-0.715$ \\
AEGIS          & Likert & $-0.706$ & $-0.413$ & $-0.468$ & $-0.737$ & $-0.635$ & $-0.680$ \\
HH-RLHF        & binary & $-0.693$ & $-0.651$ & $-0.641$ & $-0.343$ & $-0.590$ & $-0.628$ \\
HH-RLHF        & Likert & $-0.725$ & $-0.609$ & $-0.472$ & $-0.705$ & $-0.623$ & $-0.567$ \\
ToxiGen        & binary & $-0.694$ & $-0.684$ & $-0.647$ & $-0.595$ & $-0.598$ & $-0.628$ \\
ToxiGen        & Likert & $-0.672$ & $-0.718$ & $-0.864$ & $-0.786$ & $-0.800$ & $-0.621$ \\
PP (hedging)   & binary & $-0.464$ & $-0.433$ & $-0.419$ & $-0.398$ & $-0.394$ & $-0.424$ \\
PP (hedging)   & Likert & $-0.574$ & $-0.462$ & $-0.522$ & $-0.579$ & $-0.523$ & $-0.011$ \\
PP (refusal)   & binary & $-0.711$ & $-0.633$ & $-0.642$ & $-0.564$ & $-0.542$ & $-0.579$ \\
PP (refusal)   & Likert & $-0.533$ & $-0.423$ & $-0.447$ & $-0.544$ & $-0.492$ & $-0.570$ \\
MAGE           & binary & $-0.739$ & $-0.640$ & $-0.578$ & $-0.571$ & $-0.461$ & $-0.617$ \\
MAGE           & Likert & $-0.620$ & $-0.588$ & $-0.587$ & $-0.508$ & $-0.480$ & $-0.542$ \\
\bottomrule
\end{tabular}
\end{table}

Perhaps items where 9 frontier models cannot agree at L0 sit in genuinely contested or ambiguous regions of the decision boundary, and are more likely regularized
over any individual model's idiosyncrasies. A low-majority item is a signal that the underlying question is hard to label 
--- making it both a wiggle predictor and a flag for content that may be challenging to assign a confident, epistemically robust gold label in the first place.
Mechanical probes still capture a meaningful fraction of the per-item fragility (mean $|\rho| = 0.42$ for repeat, 
0.37 for position invariance), and can be a defensible single-judge fallback for evaluations without access to an LLM ensemble.

\subsection{Epistemic Fragility Beyond the Single-Shot Verdict}
\label{sec:discussion-pressure}

LLM judges occupy a unique middle ground between conventional classifiers and human raters. 
They emit a discrete verdict like a classifier, yet they can also explain it, defend it, and engage in conversation about it. 
Most LLM judges today are deployed as closed, one-shot classifiers and probably never receive turns of pushback. 
For a strictly one-shot pipeline, Mechanical Consistency is the most directly applicable part of the framework. 
However, in the framework's expansion to single-turn and multi-turn tests, we highlight two additional motivations:

\begin{enumerate}
  \item As LLMs spread to more agentic products and use cases, the oversight that a judge provides may also become more agentic. For example, in an automated moderation appeal, a safety LLM judge could issue a decision and an affected party—or an LLM acting on that party's behalf—supplies a counterargument for reconsideration.
  \item Although the paper focuses on judges, the framework bridges to deeper unsolved questions about how to measure epistemic stability. Judging makes this question more tractable because verdicts are typically discrete and, where ground-truth labels exist, changes can be classified as corrective or corrupting.
\end{enumerate}

Our results suggest that none of the LLMs we tested are consistently robust or corrective under pressure, even on canonical safety tasks. 
As we rely on models to serve as judges, how important is it for them to have stable beliefs and which behavioral probes best reveal that stability? 
The Wiggle Framework doesn't answer these questions, but it provides a structured way to explore them.
